
\section{Problem 619 (\#619)}

\subsection*{1) ROUND-2 OBJECTIVE}

Path (C): the full conjecture remains open, so I will prove it for a genuinely new infinite family that was not covered in Round~1, namely \emph{all cycles}, including the odd cycles. The target is a linear bound with a positive saving:
\[
h_4(C_n)\le \left\lceil \frac n3\right\rceil -1.
\]

This is a strict extension beyond the Round-1 bipartite result, since odd cycles are non-bipartite.

\subsection*{2) Round-1 FOUNDATION USED}

I use one Round-1 result:

\begin{enumerate}
\item If $G$ is connected bipartite on $n$ vertices, then $h_4(G)\le n/2$.
\end{enumerate}

This will only be used as context/comparison; the proof below is self-contained and addresses a new non-bipartite family.

\subsection*{3) NEW INSIGHT / TOOL (ROUND-2)}

The new construction replaces the bipartite ``join one side to a hub'' idea by a \emph{sparse dominating set of portals} on the cycle.

The portals form a dominating set $P$ of size about $n/3$, and one carefully chosen portal $c\in P$ is at cyclic distance at least $3$ from every other portal. Adding only the star edges from $c$ to $P\setminus\{c\}$ keeps the graph triangle-free, while the portal star has diameter $2$. Since every cycle vertex is within distance $1$ of a portal, the whole augmented graph has diameter at most $4$.

\subsection*{4) ATTACK PLAN (ROUND-2)}

The exact Round-1 obstacle was this: a single radius-$2$ hub is hard to realize in non-bipartite triangle-free graphs because indiscriminate hub attachments create triangles.

For cycles, the plan overcomes that obstacle by using only a \emph{very sparse} set of hub attachments.

\begin{enumerate}
\item Choose a dominating set $P$ on $C_n$ with $|P|=\lceil n/3\rceil$.
\item Choose a center $c\in P$ that is at cyclic distance at least $3$ from every other portal.
\item Add the star edges $cp$ for $p\in P\setminus\{c\}$.
\item Prove:
\begin{itemize}
\item no triangle is created, because $c$ has no common cycle-neighbor with any other portal;
\item diameter is at most $4$, because every vertex is within distance $1$ of a portal and the portal graph has diameter $2$.
\end{itemize}
\end{enumerate}

\subsection*{5) WORK (ROUND-2)}

\paragraph{Theorem 5.1.}
For every cycle $C_n$ with $n\ge 5$,
\[
h_4(C_n)\le \max\left(0,\left\lceil \frac n3\right\rceil -1\right).
\]
In particular, for odd cycles,
\[
h_4(C_n)\le \left\lceil \frac n3\right\rceil -1 < \left(1-\frac23\right)n
\]
for all sufficiently large $n$.

\paragraph{Proof.}
If $n\le 9$, then
\[
\operatorname{diam}(C_n)=\left\lfloor \frac n2\right\rfloor \le 4,
\]
so $h_4(C_n)=0$. Thus assume $n\ge 10$.

Label the cycle vertices by $0,1,\dots,n-1$ modulo $n$.

\medskip
\noindent
\textbf{Case 1: $n=3q$.}
Let
\[
P=\{0,3,6,\dots,3q-3\},
\qquad c=0.
\]
Then $|P|=q=\lceil n/3\rceil$.

\medskip
\noindent
\textbf{Case 2: $n=3q+1$.}
Let
\[
P=\{0,3,6,\dots,3q-3,\,3q-1\},
\qquad c=3.
\]
Then $|P|=q+1=\lceil n/3\rceil$.

\medskip
\noindent
\textbf{Case 3: $n=3q+2$.}
Let
\[
P=\{0,3,6,\dots,3q\},
\qquad c=3.
\]
Then $|P|=q+1=\lceil n/3\rceil$.

In every case, define $H$ by adding the star edges
\[
cp \qquad (p\in P\setminus\{c\})
\]
to the cycle.

We now verify the three required properties.

\paragraph{(i) $P$ is a dominating set of $C_n$.}
This is immediate from the construction: along the cycle, the consecutive gaps between portals have length at most $3$, so every cycle vertex lies either in $P$ or next to a portal. Equivalently, every vertex of $C_n$ is at distance at most $1$ from some portal.

\paragraph{(ii) $H$ is triangle-free.}
On a cycle, two vertices have a common neighbor if and only if they are at cyclic distance $2$.

So it is enough to check that the center $c$ is at cyclic distance at least $3$ from every other portal. This is straightforward from the case definitions:

\begin{itemize}
\item if $n=3q$, then the other portals are at distances $3,6,\dots,3q-3$ from $0$, and the wrap-around distance to the last portal is $3$;
\item if $n=3q+1$, then $c=3$ and the other portals are $0,6,9,\dots,3q-3,3q-1$; their cyclic distances from $3$ are at least $3$;
\item if $n=3q+2$, then $c=3$ and the other portals are $0,6,9,\dots,3q$; again their cyclic distances from $3$ are at least $3$.
\end{itemize}

Hence no added edge $cp$ can complete a triangle with cycle edges.

Also, because the portals are pairwise nonadjacent on the cycle, no triangle can use two added edges $cp$ and $cq$. Therefore $H$ is triangle-free.

\paragraph{(iii) $\operatorname{diam}(H)\le 4$.}
Take any two vertices $u,v$ of the cycle. By (i), choose portals $p(u),p(v)\in P$ such that
\[
\operatorname{dist}_H(u,p(u))\le 1,
\qquad
\operatorname{dist}_H(v,p(v))\le 1.
\]
Inside the portal set $P$, the added edges form a star centered at $c$, so any two portals are at distance at most $2$ in $H$. Therefore
\[
\operatorname{dist}_H(u,v)
\le
\operatorname{dist}_H(u,p(u))
+
\operatorname{dist}_H(p(u),p(v))
+
\operatorname{dist}_H(p(v),v)
\le 1+2+1=4.
\]
Thus $\operatorname{diam}(H)\le 4$.

Finally, the number of added edges is exactly
\[
|P|-1=\left\lceil \frac n3\right\rceil -1.
\]
This proves the theorem.
\hfill $\square$

\paragraph{Corollary 5.2.}
The original Erd\H{o}s--Gy\'arf\'as--Ruszink\'o question has a positive answer for the infinite family of cycles. Indeed one may take $c=2/3$ for this family.

\subsection*{6) ADVERSARIAL VERIFICATION}

\begin{enumerate}
\item The small cycles $C_5,C_6,C_7,C_8,C_9$ were separated out at the start, since they already have diameter at most $4$.
\item The only possible way to create a triangle in the construction is via one added edge and two cycle edges, or two added edges and one cycle edge. Both were checked explicitly.
\item The diameter proof uses only two facts:
\begin{itemize}
\item every vertex is within distance $1$ of a portal;
\item the portal graph is a star, hence has diameter $2$.
\end{itemize}
No hidden assumptions about parity are used.
\item The edge count is exact for the construction:
\[
|F|=|P|-1=\left\lceil \frac n3\right\rceil -1.
\]
\end{enumerate}

\subsection*{7) FINAL}

\textbf{UNRESOLVED (BUT STRICTLY ADVANCED).}

A new non-bipartite infinite family is now settled:

\[
h_4(C_n)\le \max\left(0,\left\lceil \frac n3\right\rceil -1\right)
\qquad\text{for every cycle } C_n.
\]

In particular, the conjectured linear saving holds for all odd cycles with a much stronger constant than Round~1 provided:
\[
h_4(C_n)\le \left\lceil \frac n3\right\rceil -1.
\]

The full problem for arbitrary connected triangle-free graphs remains open.

\subsection*{8) COMPLETION ESTIMATE}

COMPLETION: 56\%

\subsection*{9) REFERENCES}

\begin{itemize}
\item T. F. Bloom, \emph{Erd\H{o}s Problem \#619}, accessed 2026-03-08.
\item Round-1 notes from the previous iteration of this investigation (used only for the bipartite comparison bound listed in Section~2).
\end{itemize}

